\begin{frame}{이 예가 보여준 것: 기대값으로 값이 결정된다}
  \begin{itemize}
    \item (1) $V$ 는 각 행동을 정책 확률로 \textbf{가중평균}한 $Q$
      의 합과 \textbf{정확히 일치} --- ``동전의 양면''이 수식
      수준에서 성립.
    \item (2) $S_0$ 와 $S_1$ 의 \textbf{가치가 같다(20 $=$ 20)} ---
      ``보상이 1인 행동''과 ``보상이 3인 행동''을 1/2씩 골라서
      \textbf{기대 즉시 보상이 2}가 되는 $S_0$ 는, \textbf{항상}
      보상 2를 받는 $S_1$ 과 \textbf{기대 리턴이 같다}.
  \end{itemize}
  \vskip0.5em
  \begin{block}{상태의 가치는 ``기대(평균) 보상''으로 결정된다}
    상태의 가치는 ``매 스텝 \textbf{보장되는} 보상''이 아니라
    \textbf{정책이 정하는 \kb{기대(평균) 보상}}에 의해 결정된다.
  \end{block}
  \vskip0.4em
  \textbf{대조: 결정적 예}($R_0=1, R_1=2$)에서는
  $V(S_0) \approx 14.737 \neq 15.263 \approx V(S_1)$ 로
  \textbf{두 상태의 가치가 다르다}. 같은 ``서로 오가는 2상태''
  구조인데도, 보상(1,2)이 \textbf{고정}되어 있어서 상태마다
  \textbf{평균 보상이 달라지는} 것. 반면 이번 \textbf{확률적} 예
 에서는 $S_0$ 의 \textbf{기대} 보상이 2로 \textbf{합쳐지다 보니}
  $S_1$ 의 고정 보상 2와 같아 두 가치가 \textbf{동일}해진다.
  \vskip0.3em
  {\footnotesize 주의: 두 상태의 가치가 ``같다''는 이 특수한
  1/2-1/2 예의 \textbf{우연}이지, 일반적으로는 다르다 --- 아래
  확인 문제 (a)(c)에서 정책을 0.1-0.9로 바꾸면 다시 달라진다.}
\end{frame}
